% Generated by scripts/assemble.py; corrections belong in source/corrections.json.
\examhead{124}{中国海洋大学 2024 年数学分析试卷}{118}
% p119-b11; source PDF page 119
\par\Needspace{5\baselineskip}
\blockanchor{p119-b11}
\majorquestion{一}计算题.\par
% p119-b12; source PDF page 119
\blockanchor{p119-b12}
\question{1}求极限 $\lim_{x\to 0}\frac{1}{x}\left(\frac{1}{x} -\cos x\right).$\par
% p119-b13; source PDF page 119
\blockanchor{p119-b13}
\question{2}利用  $\frac{1}{n+1}<\ln\left(1+\frac{1}{n}\right)<\frac{1}{n}, n\in\mathbb{N}^{+}$ ，计算  $\lim_{n\to\infty}\frac{1+\frac{1}{2}+\frac{1}{3}+\cdots+\frac{1}{n}}{\ln n}$ .\par
% p119-b14; source PDF page 119
\blockanchor{p119-b14}
\question{3}已知 $\int_0^{+\infty} \frac{\sin x}{x} = \frac{\pi}{2}$, 计算 $\int_0^{+\infty} \left( \frac{\sin \alpha x}{x} \right)^2 \, \mathrm{d}x (\alpha > 0)$.\par
% p119-b15; source PDF page 119
\par\Needspace{5\baselineskip}
\blockanchor{p119-b15}
\majorquestion{二}判断题. 须加以说明.\par
% p119-b16; source PDF page 119
\blockanchor{p119-b16}
\question{1}判断正误: 设 $f(x)$ 是 $[a, b]$ 上的非负函数, 且 $\int_{a}^{b} f(x) \mathrm{d}x = 0$. 则在 $[a, b]$ 上一定存在非空开区间使得 $f$ 恒为零.\par
% p119-b17; source PDF page 119
\blockanchor{p119-b17}
\question{2}$f_{n}(x) = x^{n} - x^{3n}, 0 \leq x \leq 1$, 判断函数列 $\{f_{n}(x)\}$ 在 $[0,1]$ 上的收敛性和一致收敛性.\par
% p119-b18; source PDF page 119
\blockanchor{p119-b18}
\question{3}判断正误: 在 $[0,1]$ 上存在函数, 使得在每个有理点处连续, 但在每个无理点处不连续.\par
% p119-b19; source PDF page 119
\blockanchor{p119-b19}
\question{4}设 $a_{n} > 0$ 且级数 $\sum_{n=1}^{\infty} a_{n}$ 发散, 判断级数 $\sum_{n=1}^{\infty} \frac{a_{n}}{1 + n^{2} a_{n}}$ 的敛散性.\par
% p119-b20; source PDF page 119
\par\Needspace{5\baselineskip}
\blockanchor{p119-b20}
\majorquestion{三}证明题.\par
% p119-b21; source PDF page 119
\blockanchor{p119-b21}
\question{1}设 $f(x)$ 为 $U_{-}^{\circ}(x_0)$ 上的单调有界函数, 证明 $\lim_{x \to x_0^{-}} f(x)$ 存在.\par
% p119-b22; source PDF page 119
\blockanchor{p119-b22}
% p119-b23; source PDF page 119
\blockanchor{p119-b23}
\question{2}叙述并证明拉格朗日中值定理.\par
% p119-b24; source PDF page 119
\par\Needspace{5\baselineskip}
\par\Needspace{6\baselineskip}
\blockanchor{p119-b24}
\majorquestion{四}设 $f(x)$ 在 $[0,1]$ 上连续, 且 $f(0) = f(1)$. 证明: 对任何正整数 $n$, 必存在相应的 $\xi \in (0,1)$, 使得\par
\nopagebreak[4]
% p119-b25; source PDF page 119
\blockanchor{p119-b25}
\displayformula{f \left(\xi + \frac {1}{n}\right) = f (\xi).}
% p120-b00; source PDF page 120
\par\Needspace{5\baselineskip}
\par\Needspace{6\baselineskip}
\blockanchor{p120-b00}
\majorquestion{五}设 $f(x)$ 为 $[0, +\infty)$ 上正值单调减少函数, 且 $\int_{0}^{+\infty} f(x) \mathrm{d}x$ 收敛. 证明:\par
\nopagebreak[4]
% p120-b01; source PDF page 120
\blockanchor{p120-b01}
\displayformula{\lim _ {h \rightarrow 0 ^ {+}} h \sum_ {n = 1} ^ {\infty} f (n h) = \int_ {0} ^ {+ \infty} f (x) \mathrm{d} x.}
% p120-b02; source PDF page 120
\par\Needspace{5\baselineskip}
\par\Needspace{6\baselineskip}
\blockanchor{p120-b02}
\majorquestion{六}计算第二类曲线积分\par
\nopagebreak[4]
% p120-b03; source PDF page 120
\blockanchor{p120-b03}
\displayformula{I = \int_ {L} \frac {(x + y) \mathrm{d} x - (x - y) \mathrm{d} y}{x ^ {2} + y ^ {2}}.}
% p120-b04; source PDF page 120
\blockanchor{p120-b04}
其中 L 是从点  $A(-1,0)$  到点  $B(1,0)$  的一条在 x 轴上侧且不通过原点的光滑曲线.\par
% p120-b05; source PDF page 120
\par\Needspace{5\baselineskip}
\par\Needspace{6\baselineskip}
\blockanchor{p120-b05}
\majorquestion{七}求函数 $f(x, y, z) = x^4 + y^4 + z^4$ 在条件 $xyz = 1$ 下的极值，并判断该极值是极大值还是极小值？为什么？
\majorquestion{八}计算三重积分\par
\nopagebreak[4]
% p120-b06; source PDF page 120
\blockanchor{p120-b06}
\displayformula{\iiint_ {\Omega} z \mathrm{d} x \mathrm{d} y \mathrm{d} z.}
% p120-b07; source PDF page 120
\blockanchor{p120-b07}
其中 $\Omega$ 是球面 $x^{2} + y^{2} + z^{2} = 8$ 的上半部分与抛物面 $z = \frac{1}{2} (x^2 +y^2)$ 所围立体.\par
